\subsection{Conditional continuous-search benchmark and a negative certification result}
\label{sec:v23-source-search-pilot}
The V23 source-only continuous-search pilot used exactly the
first outer pool of BNCI2014\_002, retaining 13 held
\emph{source} people, the original covariance/tangent-LDA
models, and a fixed $[0.05,0.95]$ feasible-search range.
After five interval splits with
$\epsilon=0.02$ source-macro-AUC target tolerance,
the original frozen 27-grid source-macro selector achieved
$0.79397$ AUC. The best lambda actually evaluated by
the continuous search achieved $0.79296$ source AUC.
The initial interval enclosure left a global source-macro
AUC upper bound of $0.88773$, or residual optimality
gap $0.09477$, so \emph{no certificate was issued}.

To address excessive conservatism, we subsequently
intersected the monotone spectral bounds with an
enclosing second-order Taylor interval; this is
\emph{post-pilot mathematical refinement}, not an
independent prospective validation. With the same
five-split budget, the source-macro upper bound fell
to $0.87939$, but the remaining optimality gap of
$0.08643$ still exceeded the declared tolerance.
The empirical selected-source AUC did \emph{not}
improve, and the method honestly remained inconclusive.
The refinement improves tightness on this example,
not target-person generalisation.

\begin{table}[tb]
\centering\small
\caption{Real BNCI2014\_002 held-source continuous search,
one outer training pool, five interval splits.
Source-validation only; not a new TARGET result.}
\label{tab:v23-source-search}
\begin{tabular}{lrr}
\toprule
Source-only operation & Empirical AUC & Upper residual gap\\
\midrule
Original 27-grid macro selector & 0.79397 & ---\\
Spectral interval, five splits & 0.79296 & 0.09477\\
Spectral $+$ Taylor, five splits & 0.79296 & 0.08643\\
\bottomrule
\end{tabular}
\end{table}
